% Generated draft: requires mathematical and source review.
\documentclass[11pt]{article}
\usepackage[margin=25mm]{geometry}
\usepackage{fontspec}
\setmainfont{Latin Modern Roman}
\usepackage{amsmath,amssymb,mathtools,unicode-math}
\setmathfont{Latin Modern Math}
\usepackage{enumitem,needspace,xurl,hyperref}
\setlength{\parindent}{0pt}
\setlength{\parskip}{5pt}
\setlength{\emergencystretch}{3em}
\newcommand{\N}{\mathbb{N}}
\newcommand{\Z}{\mathbb{Z}}
\newcommand{\Q}{\mathbb{Q}}
\newcommand{\R}{\mathbb{R}}
\newcommand{\C}{\mathbb{C}}
\newcommand{\F}{\mathbb{F}}
\newcommand{\A}{\mathbb{A}}
\newcommand{\PP}{\mathbb{P}}
\newcommand{\E}{\mathbb{E}}
\newcommand{\eps}{\varepsilon}
\newcommand{\dd}{\,\mathrm{d}}
\begin{document}
\section*{UnsolvedMath ID: 4900009}
\textbf{Source problem code:} catalogue entry
\subsection{Definitions and mathematical statement}
Suppose a specified invariant ergodic probability measure of a differentiable dynamical system has finite real Lyapunov exponents \(\lambda_1\ge\cdots\ge\lambda_d\), counted with multiplicity. Put \(s_0=0\), \(s_m=\sum_{i=1}^m\lambda_i\), and
\[
 j=\max\{0\le m\le d:s_m\ge0\}.
\]
The Kaplan--Yorke number associated with this spectrum is conventionally
\[
 D_{\mathrm{KY}}=\begin{cases}
 d,&j=d,\\
 j+\dfrac{s_j}{|\lambda_{j+1}|},&j<d.
 \end{cases}
\]
In the second case \(\lambda_{j+1}<0\); if \(j=0\), the displayed value is \(0\). This definition assumes a fixed spectrum; for nonergodic measures it may vary from point to point.

The source asks for equality between this Lyapunov dimension and an ``appropriate dimension'' of an attractor or invariant measure under unspecified hypotheses. These are different targets. For instance the Hausdorff dimension of a probability measure can be defined by
\[
 \dim_H\mu=\inf\{\dim_H E:E\text{ is Borel and }\mu(E)=1\},
\]
whereas \(\dim_H A\) concerns the whole attractor \(A\). Information and box dimensions are further distinct possibilities. A measure-level proposed equality \(\dim_H\mu=D_{\mathrm{KY}}(\mu)\) requires a specified class of systems and measures; a set-level statement also needs a rule selecting or combining spectra of measures on the attractor. The supplied excerpt gives neither.

The research summary reports upper-bound results in broad classes and attributes equality results for SRB and sufficiently regular or absolutely continuous invariant measures to Ledrappier--Young, without exact hypotheses or dimension definitions. Here SRB refers to the usual Sinai--Ruelle--Bowen measure terminology, involving absolute continuity of conditional measures in unstable directions when those are defined. The summary also explicitly says equality can fail for non-SRB measures. Therefore an unrestricted equality for arbitrary invariant measures cannot be treated simultaneously as the intended universally valid conjecture. The admissible class and the exact equality target remain unresolved in this draft; the partial-progress label is a source attribution, not a verified theorem or current status assessment.
\subsection{Short English statement}
When does the Kaplan--Yorke number computed from Lyapunov exponents equal a geometric dimension of an attractor or invariant measure? The record does not choose the dimension or hypotheses and itself reports failures for some measures, so a universal equality cannot be substituted for a precise restricted conjecture.
\subsection{Sources}
\begin{itemize}
\item Local registry: \url{data/unsolvedmath/problems.json}, permanent UnsolvedMath ID \(4900009\), source code catalogue entry; supplied statement, background, and research summary, updated 2026-08-06. Proposed by James Kaplan and James Yorke.
\item Wikipedia contributors, \emph{List of unsolved problems in mathematics}, Dynamical systems, current source bullet 9: \url{https://en.wikipedia.org/wiki/List_of_unsolved_problems_in_mathematics#Dynamical_systems}. HTML extraction accessed 2026-07-29. The record explicitly cautions that Wikipedia is not an authoritative resolution source.
\item Ledrappier--Young is named in the supplied summary, but no original-paper citation or exact equality theorem is supplied. No authoritative original formulation of the conjecture is included.
\end{itemize}
\end{document}
